% -------------------------------------------------------------
% ABSTRACT
% -------------------------------------------------------------
\thispagestyle{plain}
\hspace*{1cm} Efficient heat production planning is a significant challenge in modern district heating systems. In the city of Heatington, this planning is still performed manually, which makes the process time-consuming and prone to inefficiencies while supplying heat to approximately 1600 buildings. This project - HeatItOn, represents the development of a web application, providing an automated heat production platform. The project was developed using .NET C\# and a MySQL database. API clients were used to provide communication between the backend and frontend components. The frontend was developed using Avalonia \cite{avalonia}, providing a user-friendly interface with separate tabs for scenario management, source data input, and results visualization. The application consists of several modules: the Asset Manager responsible for managing assets; the Source Data Manager responsible for processing incoming data; the Optimizer responsible for the main heat production optimization logic; the Result Data Manager responsible for processing output data and the Data Visualization module responsible for presenting data through graphs and data grids. The development system enables users to input source data, select assets for optimization, and analyze the generated results through data grids and graphical visualizations. The developed system successfully automated the heat production planning workflow and generated optimized heat schedules based on user-defined scenarios. The results demonstrate that automated planning can reduce manual effort and provide clearer insights into heat production performance. Overall, the project shows that such a tool can support more efficient and data-driven decision-making in district heating operations.
